%% 第7巻 巻末（こたえ・記号カード・定数と単位・公式・チェック表・つぎの巻へ・さくいん・奥付）（第14回編集委員会）
%% 第4巻・第6巻の巻末と同じ形。記号カードの一覧・書くときのちゅうい・単位と数・公式・表計算と Python・つぎの巻へを、第7巻の本文（7-0〜7-11）に合わせて書いた
%% 第7巻はルビなし（\R は使わない）。字の色は bkTheme でなく \bkTT。つぎの巻へのカードは第8巻の色（bkVol8・bkDeep8）
%% 数と式は本文と同じ（Python で確かめた）。行列は細字、ベクトルは矢印、単位行列は I（第14回の決定）
\makeatletter
%% ---------------- こたえ ----------------
\bkAnswerPages

%% ---------------- 記号カード ----------------
%% 前の巻の記号で、この巻でもよく使うもの（章の番号は、はじめて出た巻の章。この巻の章の \bkkigocard は自動でうしろに集まる）
%% 7-7 のカードは記号を \fontsize{14}{16} で大きく見せている。一覧の記号の列（18 mm）には入らないので、
%% 一覧の中だけ、14pt の指定をそのときの字の大きさ（\footnotesize）に読みかえる（章の本文のカードはそのまま。第6巻と同じ方法）
\let\bk@origfontsize\fontsize
\def\bk@fsfourteen{14}%
\protected\def\fontsize#1#2{\def\bk@fsarg{#1}\ifx\bk@fsarg\bk@fsfourteen\expandafter\bk@origfontsize\expandafter{\f@size}{\f@baselineskip}\else\bk@origfontsize{#1}{#2}\fi}%
\bkKigoPages{%
\bkkigorow{$\Delta$}{デルタ}{変化量。この巻の $\Delta t$ は、オイラー法の1歩の時間（7-10）}{2-3}%
\bkkigorow{$\theta$}{シータ}{角の大きさ。回転行列 $R(\theta)$ の角、関節の角 $\theta_1$・$\theta_2$。7-7 ではピッチの角}{4-2}%
\bkkigorow{\shortstack{$\cos\theta$\\[2pt]$\sin\theta$}}{コサイン・シータ、サイン・シータ}{回転行列の成分。表計算や Python では、度をラジアンに直してから計算する（7-6）}{4-6}%
\bkkigorow{$\omega$}{オメガ}{角速度〔rad/s〕。関節の角速度 $\dot\theta_1$、$\dot\theta_2$（7-2）や、台車の角振動数（7-8）}{4-8}%
\bkkigorow{$\vec a$}{ベクトル・エー}{向きと大きさをもつ量。この巻では、成分をたてに書く（7-1）}{4-10}%
\bkkigorow{$\vec e_x$、$\vec e_y$}{イー・エックス、イー・ワイ}{$x$・$y$ の向きの大きさ1の矢印（物差しの矢印）。行列の1列めは $\vec e_x$ の行き先（7-2）}{4-11}%
\bkkigorow{$\vec a\cdot\vec b$}{エー・ドット・ビー}{内積。行列の積の成分は、行と列の内積（7-3）}{4-12}%
\bkkigorow{$\vec a\times\vec b$}{エー・クロス・ビー}{外積。箱の体積 $(\vec a\times\vec b)\cdot\vec c$ は、3×3 の行列式（7-4）}{4-14}%
\bkkigorow{$\dot x$、$\ddot x$}{エックス・ドット、エックス・ツードット}{時間で1回・2回微分。$\dot{\vec x}$ は成分をそれぞれ微分したもの（7-8）}{5-5}%
\bkkigorow{$\Sigma$}{シグマ}{全部たす。積の成分 $\sum_k a_{ik}b_{kj}$（7-1）}{5-10}%
\bkkigorow{$y_n$}{ワイ・エヌ}{$n$ 歩めの値。この巻の $\vec x_n$ は、そのベクトル版（7-10）}{6-2}%
\bkkigorow{$j$}{ジェイ}{虚数単位。$R_{90}R_{90}=-I$ は、$j^2=-1$ と同じ（7-6）}{6-8}%
\bkkigorow{$\bar z$}{ゼット・バー}{共役な複素数。共役な固有値の組では $\lambda\bar\lambda=|\lambda|^2$（7-10）}{6-8}%
\bkkigorow{$|z|$}{ぜったいち・ゼット}{複素数の大きさ。固有値の大きさ $|\lambda|$ が、くり返しのふえ方を決める（7-10）}{6-9}%
\bkkigorow{$e^{j\theta}$}{イー・ジェイ・シータ}{$\cos\theta+j\sin\theta$。かけると $\theta$ 回る。回転行列と同じ回転（7-6・7-10）}{6-10}%
\bkkigorow{$j\omega$}{ジェイ・オメガ}{コイルの箱の $j\omega L$、コンデンサーの箱の $j\omega C$（7-3）}{6-10}%
\bkkigorow{$\cosh x$}{ハイパボリック・コサイン}{送電線の箱の行列の成分 $\cosh\gamma\ell$、$\sinh\gamma\ell$（7-11）}{6-11}%
\bkkigorow{$s$}{エス}{特性方程式の文字。状態の行列の固有値 $\lambda$ と同じもの（7-8）}{6-12}%
}
\let\fontsize\bk@origfontsize

\def\bkbkhead#1{\par\addvspace{4mm}\noindent{\sffamily\bfseries\large\color{\bkTT}#1}\par\nobreak\vspace{2mm}}
%% 記号カードのつづき：手で書くとき・読むときのちゅうい（各章の「まちがえやすいところ」から）
\Needspace{8\baselineskip}
\bkbkhead{記号を書くときのちゅうい}
{\small 手で書くと、ほかの記号とまぎれやすいものと、読みまちがえやすい書き方です。行列では、書く順番がそのまま意味になります。}\par\vspace{2mm}
{\small
\begin{longtblr}{colspec={Q[c,m,wd=21mm] Q[c,m,wd=22mm] Q[l,m,wd=83mm]}, rowhead=1, row{1}={bg=bkVol!14, font=\sffamily\bfseries\footnotesize}, hlines={0.4pt,bkLine}, rowsep=0.8pt}
記号 & まぎれやすい & 書くとき・読むときのコツ \\
$a_{12}$ & $a_{21}$ & 1行め2列め。行の番号が先（表計算の番地 B3 は列が先）（7-1） \\
$A\vec x$ & $\vec x A$ & ベクトルは行列の右に、たてに書く。機械（行列）に右から入れる（7-1・7-2） \\
$AB$ & $BA$ & ふつう $AB\neq BA$。$AB\vec x$ は右の $B$ が先（先にする変換を右に書く）（7-3） \\
$A^2$ & 成分の2乗 & $A^2=AA$。主役の行列なら1行めは $(5,\ 4)$。成分の2乗 $(4,\ 1)$ ではない（7-3・7-10） \\
$A^{\mathsf T}$ & $A^T$ & 立てた字の T。「$T$ 乗」ではない。行と列を入れかえる（7-1） \\
$|A|$ & 絶対値 & 行列式 $\det A$ の、専門書の書き方。マイナスにもなる（7-4） \\
$\det(2A)$ & $2\det A$ & 2×2 では $2^2\det A$、3×3 では $2^3\det A$。$\det(A+B)$ にはきまりがない（7-4） \\
$A^{-1}$ & $\dfrac{1}{A}$ & 行列のわり算はない。$A\vec x=\vec b$ なら $\vec x=A^{-1}\vec b$（左からかける）（7-5） \\
$(AB)^{-1}$ & $A^{-1}B^{-1}$ & $B^{-1}A^{-1}$。あとからしたことを先にもどす（くつ下とくつ）（7-5） \\
$R(\theta)$ & $R(-\theta)$ & 1行めは $(\cos\theta,\ -\sin\theta)$。90°を入れ、$(1,\ 0)$ が $(0,\ 1)$ へ行くか確かめる（7-6） \\
$R(\theta)^{-1}$ & 成分の逆数 & $R(\theta)^{-1}=R(\theta)^{\mathsf T}=R(-\theta)$（7-6） \\
$R_y(\theta)$ & $R_x$、$R_z$ の形 & 1行めは $(\cos\theta,\ 0,\ \sin\theta)$。$\sin$ の符号が逆に見える（7-7） \\
$(x,\ y,\ 1)$ & $(x,\ y,\ 0)$ & 同次座標。位置には1、向き（速度など）には0をつける（7-7） \\
$A-\lambda I$ & $A-\lambda$ & $\lambda$ は対角成分からだけひく。ほかの成分はそのまま（7-8） \\
$\vec v$ & $\vec 0$ & 固有ベクトルに $\vec 0$ はふくめない。何倍しても（向きが逆でも）同じ固有ベクトル（7-8） \\
$P^{-1}AP$ & $PAP^{-1}$ & $D=P^{-1}AP$。$P$ の列の順と、$D$ の対角成分の順をそろえる（7-9） \\
$|\lambda|$ & $\lambda$ & ふえるか消えるかは大きさで決まる。$|-1.5|=1.5>1$ はふえる。$|1+0.1j|=\sqrt{1.01}$（7-10） \\
$I+A\,\Delta t$ & $A\,\Delta t$ & オイラー法の1歩の行列。今の値 $I\vec x_n$ をたすのを忘れない（7-10） \\
\end{longtblr}\par}

%% 記号カードのつづき：ギリシャ文字（この巻での意味に合わせた）
\par\noindent\begin{minipage}{\linewidth}
\bkbkhead{ギリシャ文字の読み方}
{\small この巻では、固有値 $\lambda$、姿勢の角 $\phi$・$\theta$・$\psi$ など、ギリシャ文字がたくさん出てきます。大文字・小文字・読み方と、この巻での意味です（かっこの中は章）。}\par\vspace{2mm}
\begin{bkcenter}
\footnotesize
\begin{tblr}{colspec={Q[c,m,wd=5mm] Q[c,m,wd=5mm] Q[l,m,wd=15mm] Q[l,m,wd=37mm] Q[c,m,wd=5mm] Q[c,m,wd=5mm] Q[l,m,wd=15mm] Q[l,m,wd=37mm]}, row{1}={bg=bkVol!14, font=\sffamily\bfseries\footnotesize}, hlines={0.4pt,bkLine}, vline{5}={0.8pt,bkVol}, rowsep=0.6pt, colsep=2pt}
大 & 小 & 読み方 & この巻での意味 & 大 & 小 & 読み方 & この巻での意味 \\
$\mathrm{A}$ & $\alpha$ & アルファ & 回す角（7-6） & $\mathrm{N}$ & $\nu$ & ニュー & ― \\
$\mathrm{B}$ & $\beta$ & ベータ & 回す角（7-6） & $\Xi$ & $\xi$ & クシー & ― \\
$\Gamma$ & $\gamma$ & ガンマ & 送電線の $\gamma$（7-11） & $\mathrm{O}$ & $o$ & オミクロン & ― \\
$\Delta$ & $\delta$ & デルタ & $\Delta t$、$\Delta x$（7-10・7-11） & $\Pi$ & $\pi$ & パイ & 円周率 \\
$\mathrm{E}$ & $\varepsilon$ & イプシロン & ― & $\mathrm{P}$ & $\rho$ & ロー & ― \\
$\mathrm{Z}$ & $\zeta$ & ゼータ & （減衰比。6-7） & $\Sigma$ & $\sigma$ & シグマ & $\Sigma$ は和、$\sigma$ は応力（7-2） \\
$\mathrm{H}$ & $\eta$ & イータ & ― & $\mathrm{T}$ & $\tau$ & タウ & （時定数。6-4） \\
$\Theta$ & $\theta$ & シータ & 回す角・関節の角、ピッチ（7-7） & $\Upsilon$ & $\upsilon$ & ウプシロン & ― \\
$\mathrm{I}$ & $\iota$ & イオタ & ― & $\Phi$ & $\phi$ & ファイ & ロール（7-7） \\
$\mathrm{K}$ & $\kappa$ & カッパ & ― & $\mathrm{X}$ & $\chi$ & カイ & ― \\
$\Lambda$ & $\lambda$ & ラムダ & 固有値（7-8） & $\Psi$ & $\psi$ & プサイ & ヨー（7-7） \\
$\mathrm{M}$ & $\mu$ & ミュー & マイクロ（\unit{\micro F}） & $\Omega$ & $\omega$ & オメガ & 角速度。$\Omega$ はオーム \\
\end{tblr}
\end{bkcenter}
{\footnotesize 大文字が英語のアルファベットと同じ形の文字（A、B、E など）は、記号としてはほとんど使いません。ロールの $\phi$ は、6-6 の初期位相の $\phi$ とは別のものです。$\phi$ は $\varphi$ と書くこともあります（4-15 の球座標の角は $\varphi$）。応力の行列は、物理の本では成分の $\sigma_x$ と区別するために太字 $\symbfit{\sigma}$ で書きます（7-2）。}
\end{minipage}


%% ---------------- 定数・単位・公式 ----------------
\bkbacktitle{よく使う定数・単位・公式}
\bksubsection{1}{この巻で使う単位}
{\small
\begin{longtblr}{colspec={Q[l,m,wd=25mm] Q[l,m,wd=32mm] Q[l,m,wd=53mm] Q[c,m,wd=12mm]}, rowhead=1, row{1}={bg=bkVol!14, font=\sffamily\bfseries\footnotesize}, hlines={0.4pt,bkLine}, rowsep=1.4pt}
量 & 単位（読み方） & ほかの単位との関係 & 章 \\
長さ & m、cm、mm & $\qty{1}{m}=\qty{100}{cm}=\qty{1000}{mm}$ & 7-2 \\
角 & °（度）、rad（ラジアン） & $\pi$\,rad＝180°。表計算・Python の cos はラジアン & 7-6 \\
角速度 & rad/s & 関節の回る速さ $\dot\theta$。1秒に回る角 & 7-2 \\
速度 & m/s & 手の速度 $\dot{\vec p}=J\dot{\vec\theta}$ & 7-2 \\
力・重さ & N（ニュートン） & 3つのばねの球1個の重さ \qty{0.98}{N} & {7-1\\7-5} \\
ばね定数 & N/m & \qty{1}{m}のばすのに要る力 & {7-5\\7-8} \\
応力 & MPa（メガパスカル） & $\qty{1}{MPa}=\qty{1}{N/mm^2}$ & 7-2 \\
電圧・電流 & V（ボルト）、A（アンペア） & 箱の入口の組 $(V_1,\ I_1)$、出口の組 $(V_2,\ I_2)$ & 7-3 \\
抵抗 & $\Omega$（オーム） & $\qty{1}{\ohm}=\qty{1}{V/A}$。箱の行列の右上の単位 & 7-3 \\
コンダクタンス & S（ジーメンス） & $\qty{1}{S}=\qty{1}{A/V}$（$\frac{1}{\Omega}$）。箱の行列の左下、$j\omega C$ の単位 & 7-3 \\
インダクタンス・静電容量 & H（ヘンリー）、F（ファラド） & $\qty{1}{mH}=\qty{0.001}{H}$、\ $\qty{1}{\micro F}=\qty{0.000001}{F}$ & 7-3 \\
角周波数 & rad/s & $\omega=2\pi f$（$f$ は Hz） & 7-3 \\
\end{longtblr}\par}

\bksubsection{2}{よく使う数}

\begin{bkcenter}
\small
\begin{tblr}{colspec={Q[l,m,wd=36mm] Q[l,m,wd=76mm] Q[c,m,wd=12mm]}, row{1}={bg=bkVol!14, font=\sffamily\bfseries\footnotesize}, hlines={0.4pt,bkLine}, rowsep=1pt}
名前 & 値 & 章 \\
$\dfrac{1}{\sqrt2}$ & 0.70710678…（numpy の固有ベクトル） & 7-8 \\
$\cosh1$、\ $\sinh1$ & 1.543、\ 1.175（送電線の箱の行列） & 7-11 \\
重力加速度 $g$ & \qty{9.8}{m/s^2}（\qty{1.0}{kg} の荷物の重力 \qty{9.8}{N}） & 7-6 \\
掃き出し法の手間（$n$ 文字） & 約 $\dfrac{n^3}{3}$ 回のかけ算（$n=100$ で約33万回） & 7-5 \\
\end{tblr}
\end{bkcenter}
\vspace{2mm}
\noindent\begin{minipage}{\linewidth}
\noindent{\sffamily\bfseries\small\color{\bkTT}回転行列に入る三角比の値}（4-6・7-6）\par\vspace{1mm}
\begin{bkcenter}
\begin{tblr}{colspec={Q[l,m,wd=30mm] *{5}{Q[c,m,wd=17mm]}}, row{1}={bg=bkVol!14, font=\sffamily\bfseries\footnotesize}, hlines={0.4pt,bkLine}, rowsep=1.4pt, cells={font=\small}}
角 $\theta$ & 0° & 30° & 45° & 60° & 90° \\
ラジアン & 0 & $\dfrac{\pi}{6}$ & $\dfrac{\pi}{4}$ & $\dfrac{\pi}{3}$ & $\dfrac{\pi}{2}$ \\
$\cos\theta$ & 1 & $\dfrac{\sqrt3}{2}$ & $\dfrac{1}{\sqrt2}$ & $\dfrac{1}{2}$ & 0 \\
（小数） & 1 & 0.866 & 0.707 & 0.5 & 0 \\
$\sin\theta$ & 0 & $\dfrac{1}{2}$ & $\dfrac{1}{\sqrt2}$ & $\dfrac{\sqrt3}{2}$ & 1 \\
（小数） & 0 & 0.5 & 0.707 & 0.866 & 1 \\
\end{tblr}
\end{bkcenter}
{\footnotesize 行列の成分は、小数第3位まで書きます。途中で丸めたときの小さなちがいは、丸めのためです（7-2・7-6）。}
\end{minipage}\par\vspace{1mm}

\Needspace{8\baselineskip}
\bkbkhead{この巻でくり返し使う場面と数}
{\small 同じ場面は、どの章でも同じ数で計算しています。}\par\vspace{1mm}
{\small
\begin{longtblr}{colspec={Q[l,m,wd=20mm] Q[l,m,wd=93mm] Q[c,m,wd=12mm]}, rowhead=1, row{1}={bg=bkVol!14, font=\sffamily\bfseries\footnotesize}, hlines={0.4pt,bkLine}, rowsep=1.6pt}
場面 & 行列と数 & 章 \\
くだもの屋 & 個数 $N=\begin{pmatrix}2&3\\2&1\end{pmatrix}$、値段（りんご120円・みかん60円）で合計420円・300円。お店B（100円・80円）もならべると $N\begin{pmatrix}120&100\\60&80\end{pmatrix}=\begin{pmatrix}420&440\\300&280\end{pmatrix}$ & 7-1 \\
3-6 の回路 & $\begin{pmatrix}4&2\\2&4\end{pmatrix}\begin{pmatrix}I_1\\I_2\end{pmatrix}=\begin{pmatrix}9\\6\end{pmatrix}$。行列式12、$I_1=\qty{2}{A}$、$I_2=\qty{0.5}{A}$ & {7-1\\7-5} \\
3つの文字 & $x+y+z=9$、$x-y=1$、$y-z=1$ を掃き出し法で：$x=4$、$y=3$、$z=2$ & 7-5 \\
4-11 の看板 & 30°と60°のひもで \qty{100}{N}。行列式 $-1$、$T_1=\qty{50}{N}$、$T_2=\qty{86.6}{N}$。ひもが水平に近いと行列式は0に近く、力が大きい & {7-1\\7-4\\7-5} \\
2本の腕 & $L_1=\qty{0.30}{m}$、$L_2=\qty{0.25}{m}$、$\theta_1=30°$、$\theta_2=45°$。手 $(0.325,\ 0.391)$。$J=\begin{pmatrix}-0.391&-0.241\\0.325&0.065\end{pmatrix}$、$\det J=L_1L_2\sin\theta_2=0.0530$ & {7-2\\7-4\\7-6\\7-7} \\
電気の箱 & 直列 \qty{2}{\ohm} $\begin{pmatrix}1&2\\0&1\end{pmatrix}$、並列 \qty{2}{\ohm} $\begin{pmatrix}1&0\\0.5&1\end{pmatrix}$。直列が入口側で $\begin{pmatrix}2&2\\0.5&1\end{pmatrix}$、逆の順で $\begin{pmatrix}1&2\\0.5&2\end{pmatrix}$。行列式はどちらも1 & {7-3\\7-4} \\
3つのばね & $k=\qty{100}{N/m}$、球の重さ \qty{0.98}{N}。のび $u_1=\qty{0.0294}{m}$、$u_2=\qty{0.0490}{m}$、$u_3=\qty{0.0588}{m}$ & {7-1\\7-5} \\
2つの台車 & $m=\qty{1}{kg}$、$k=\qty{100}{N/m}$、$K=\qty{10}{N/m}$。$B=\begin{pmatrix}-110&10\\10&-110\end{pmatrix}$、固有値 $-100$・$-120$、$\omega=10$・10.95\,rad/s。固有ベクトル $(1,\ 1)$（いっしょ）・$(1,\ -1)$（反対） & {7-1\\7-8\\7-9} \\
2つの水そう & $H=\begin{pmatrix}-0.1&0.1\\0.1&-0.1\end{pmatrix}$、固有値0（合計は一定）・$-0.2$（差は $e^{-0.2t}$）。$\Delta t=1$ 分の1歩の行列の固有値は1と0.8 & {7-8\\7-10} \\
3-10 の式と、ばね & $\begin{pmatrix}0&1\\-2&-3\end{pmatrix}$ の固有値 $-1$、$-2$。$\begin{pmatrix}0&1\\-5&-2\end{pmatrix}$ は $-1\pm2j$。\qty{0.5}{kg}・\qty{8}{N/m}・$c=0.4$ の $\begin{pmatrix}0&1\\-16&-0.8\end{pmatrix}$ は $-0.4\pm3.980j$ & 7-8 \\
応力 & $\begin{pmatrix}80&30\\30&-20\end{pmatrix}$〔MPa〕。固有値（主応力）88.31・$-28.31$、固有ベクトル $(1,\ 0.277)$・$(1,\ -3.61)$ & {7-2\\7-8} \\
天気 & $M=\begin{pmatrix}0.8&0.4\\0.2&0.6\end{pmatrix}$（どの列もたすと1）。固有値1・0.4、落ち着く先は晴れ $\dfrac{2}{3}$、雨 $\dfrac{1}{3}$ & 7-10 \\
冷めるお湯 & 部屋20\,℃、はじめ80\,℃。差は1歩で $1-0.05\,\Delta t$ 倍。$\Delta t=1$ 分で10分後 55.92\,℃。$\Delta t=30$ 分で $-0.5$ 倍（ゆれる）。こわれないのは $\Delta t<40$ 分 & 7-10 \\
ばねのオイラー法 & $\omega=1$。ふつうの方法の1歩の行列 $\begin{pmatrix}1&\Delta t\\-\Delta t&1\end{pmatrix}$、固有値 $1\pm j\Delta t$。1周で1.37倍（$\Delta t=0.1$）・1.03倍（0.01）。速度を先に進める方法は行列式1、$|\lambda|=1$ & 7-10 \\
最小二乗法 & 3-16 のデータ7点。正規方程式 $\begin{pmatrix}91&21\\21&7\end{pmatrix}\begin{pmatrix}a\\b\end{pmatrix}=\begin{pmatrix}136.2\\31.5\end{pmatrix}$、$y=1.489x+0.032$ & 7-5 \\
\end{longtblr}\par}


\Needspace{14\baselineskip}
\bkbkhead{主役の行列の6つの顔}
{\small この巻の主役の行列 $A=\begin{pmatrix}2&1\\1&2\end{pmatrix}$ を、章ごとにちがう目で見ました（7-0 の表の答え合わせ）。}\par\vspace{1mm}
\begin{bkcenter}
\small
\begin{tblr}{colspec={Q[c,m,wd=11mm] Q[l,m,wd=30mm] Q[l,m,wd=83mm]}, row{1}={bg=bkVol!14, font=\sffamily\bfseries\footnotesize}, hlines={0.4pt,bkLine}, rowsep=1.6pt}
章 & 見る目 & 計算と、わかったこと \\
7-2 & 形を変える機械 & 家のかど $(2,\ 0)\to(4,\ 2)$、$(2,\ 2)\to(6,\ 6)$、$(1,\ 3)\to(5,\ 7)$。1列め $(2,\ 1)$ は $\vec e_x$ の行き先、2列め $(1,\ 2)$ は $\vec e_y$ の行き先 \\
7-4 & 面積が何倍になるか & $\det A=2\times2-1\times1=3$。正方形1つが面積3の平行四辺形に \\
7-5 & 元にもどす機械 & $A^{-1}=\dfrac{1}{3}\begin{pmatrix}2&-1\\-1&2\end{pmatrix}$。かけると $I$ \\
7-8 & 向きの変わらない矢印 & 固有値3と1。固有ベクトル $(1,\ 1)$ と $(1,\ -1)$（直交する） \\
7-9 & のびるだけに分ける & $P=\begin{pmatrix}1&1\\1&-1\end{pmatrix}$、\ $P^{-1}AP=\begin{pmatrix}3&0\\0&1\end{pmatrix}=D$。$q_1=\dfrac{x+y}{2}$ が3倍、$q_2=\dfrac{x-y}{2}$ は1倍 \\
7-10 & 何回もかける & $A^2=\begin{pmatrix}5&4\\4&5\end{pmatrix}$、\ $A^5=\begin{pmatrix}122&121\\121&122\end{pmatrix}$（成分は $\dfrac{3^5\pm1}{2}$）。点は $(1,\ 1)$ の向きにそろう \\
\end{tblr}
\end{bkcenter}


\bksubsection{3}{行列の計算の公式}
{\small
\begin{longtblr}{colspec={Q[l,m,wd=26mm] Q[l,m,wd=86mm] Q[c,m,wd=12mm]}, rowhead=1, row{1}={bg=bkVol!14, font=\sffamily\bfseries\footnotesize}, hlines={0.4pt,bkLine}, rowsep=1.6pt}
公式・きまり & 式 & 章 \\
行列と成分 & $m\times n$ 行列は $m$ 行 $n$ 列。$a_{ij}$ は $i$ 行め $j$ 列めの成分 & 7-1 \\
和・実数倍 & 成分ごと。$k\begin{pmatrix}a&b\\c&d\end{pmatrix}=\begin{pmatrix}ka&kb\\kc&kd\end{pmatrix}$。たせるのは同じ大きさどうし & 7-1 \\
行列×ベクトル & $\begin{pmatrix}a&b\\c&d\end{pmatrix}\begin{pmatrix}x\\y\end{pmatrix}=\begin{pmatrix}ax+by\\cx+dy\end{pmatrix}$（行ごとに、ベクトルとの内積） & {7-1\\7-2} \\
列の意味 & 1列めは $\vec e_x$ の行き先、2列めは $\vec e_y$ の行き先 & 7-2 \\
線形 & $A(\vec u+\vec v)=A\vec u+A\vec v$、\ $A(k\vec v)=kA\vec v$。原点は動かない（ずらす変換は 7-7 の同次座標で） & 7-2 \\
行列×行列 & $\begin{pmatrix}a&b\\c&d\end{pmatrix}\begin{pmatrix}p&q\\r&s\end{pmatrix}=\begin{pmatrix}ap+br&aq+bs\\cp+dr&cq+ds\end{pmatrix}$。$(i,\ j)$ 成分は、左の $i$ 行めと右の $j$ 列めの内積 & {7-1\\7-3} \\
大きさ & （$m$ 行 $n$ 列）×（$n$ 行 $p$ 列）＝（$m$ 行 $p$ 列）。内側の数がちがうと、かけられない & {7-1\\7-3} \\
積の意味 & $AB\vec x=A(B\vec x)$。右の $B$ が先、左の $A$ があと & 7-3 \\
積のきまり & ふつう $AB\neq BA$。$(AB)C=A(BC)$。$AI=IA=A$ & 7-3 \\
転置 & $A^{\mathsf T}$ は行と列を入れかえる。$A^{\mathsf T}=A$ なら対称。$(AB)^{\mathsf T}=B^{\mathsf T}A^{\mathsf T}$ & {7-1\\7-3} \\
単位行列・零行列 & $I=\begin{pmatrix}1&0\\0&1\end{pmatrix}$（何もしない変換）、$O$ はすべて0 & {7-1\\7-3} \\
連立方程式 & $2x+3y=420$、$2x+y=300$ なら $\begin{pmatrix}2&3\\2&1\end{pmatrix}\begin{pmatrix}x\\y\end{pmatrix}=\begin{pmatrix}420\\300\end{pmatrix}$。$A\vec x=\vec b$（$A$ は係数行列） & 7-1 \\
\end{longtblr}\par}

\vspace{2mm}
\noindent\begin{minipage}{\linewidth}
\noindent{\sffamily\bfseries\small\color{\bkTT}電気の箱の行列（四端子定数）}（7-3・7-4・7-11）\par\vspace{1mm}
{\small 箱の行列は「出口の組 $(V_2,\ I_2)$ を入れると、入口の組 $(V_1,\ I_1)$ が出てくる」向きです。箱をつなぐときは、入口側の箱を左から順にかけます。電気の本では、この4つの数を $A$・$B$・$C$・$D$ とよびます。}\par\vspace{1mm}
\begin{bkcenter}
\begin{tblr}{colspec={Q[c,m,wd=28mm] Q[c,m,wd=28mm] Q[c,m,wd=30mm] Q[c,m,wd=36mm]}, row{1}={bg=bkVol!14, font=\sffamily\bfseries\footnotesize}, hlines={0.4pt,bkLine}, vlines={0.4pt,bkLine}, rowsep=2pt, colsep=3pt, cells={font=\small}}
直列の抵抗 $R$ & 並列の抵抗 $R$ & 直列のコイル $L$ & 並列のコンデンサー $C$ \\
$\begin{pmatrix}1&R\\0&1\end{pmatrix}$ & $\begin{pmatrix}1&0\\\dfrac{1}{R}&1\end{pmatrix}$ & $\begin{pmatrix}1&j\omega L\\0&1\end{pmatrix}$ & $\begin{pmatrix}1&0\\j\omega C&1\end{pmatrix}$ \\
\end{tblr}
\end{bkcenter}
{\footnotesize どの箱も、行列式は1です。つないだ箱の行列式も $1\times1=1$ です（7-4）。}\par\vspace{1mm}
\begin{bkcenter}
\begin{tblr}{colspec={Q[l,m,wd=48mm] Q[c,m,wd=74mm]}, column{1}={font=\sffamily\bfseries\footnotesize}, hlines={0.4pt,bkLine}, vlines={0.4pt,bkLine}, rowsep=2pt, cells={font=\small}}
{コイルの箱のあとに\\コンデンサーの箱（7-3）} & $\begin{pmatrix}1-\omega^2LC&j\omega L\\j\omega C&1\end{pmatrix}$ \\
{長い送電線（7-11）} & {$\begin{pmatrix}\cosh\gamma\ell&Z_0\sinh\gamma\ell\\[2pt]\dfrac{1}{Z_0}\sinh\gamma\ell&\cosh\gamma\ell\end{pmatrix}$\\[2pt]$\gamma=\sqrt{zy}$、\ $Z_0=\sqrt{\dfrac{z}{y}}$} \\
\end{tblr}
\end{bkcenter}
\end{minipage}

\bksubsection{4}{行列式と逆行列の公式}
{\small
\begin{longtblr}{colspec={Q[l,m,wd=26mm] Q[l,m,wd=86mm] Q[c,m,wd=12mm]}, rowhead=1, row{1}={bg=bkVol!14, font=\sffamily\bfseries\footnotesize}, hlines={0.4pt,bkLine}, rowsep=1.6pt}
公式・きまり & 式 & 章 \\
2×2 の行列式 & $\det\begin{pmatrix}a&b\\c&d\end{pmatrix}=ad-bc$（ななめにかけてひく） & 7-4 \\
意味 & 面積の倍率。マイナスは裏返し。0なら2つの列が平行で、平らにつぶれる & 7-4 \\
積と数倍 & $\det(AB)=\det A\times\det B$。2×2 では $\det(kA)=k^2\det A$ & 7-4 \\
3×3 の行列式 & 3本の矢印 $\vec a$、$\vec b$、$\vec c$ を列に並べると、$a_1b_2c_3+b_1c_2a_3+c_1a_2b_3-c_1b_2a_3-a_1c_2b_3-b_1a_2c_3$（サラスの方法。3×3 だけ）。体積の倍率 & 7-4 \\
連立方程式と行列式 & $\det A\neq0$ なら答えは1つ。$\det A=0$ なら、答えがないか、無数にある & {7-4\\7-5} \\
ロボットの腕 & $\det J=L_1L_2\sin\theta_2$。0になる姿勢（腕がのびきる）が特異姿勢 & 7-4 \\
逆行列 & $\begin{pmatrix}a&b\\c&d\end{pmatrix}^{-1}=\dfrac{1}{ad-bc}\begin{pmatrix}d&-b\\-c&a\end{pmatrix}$。$A^{-1}A=AA^{-1}=I$ & 7-5 \\
逆行列の作り方 & {①左上と右下を入れかえる\\②右上と左下にマイナスをつける\\③行列式でわる\\作ったら、かけて $I$ になるか確かめる} & 7-5 \\
連立方程式の答え & $A\vec x=\vec b$ の答えは $\vec x=A^{-1}\vec b$（左からかける） & 7-5 \\
順番 & $(AB)^{-1}=B^{-1}A^{-1}$ & 7-5 \\
掃き出し法 & {係数と右辺を並べた表に、3つの操作をして左を $I$ にする。右の列が答え\\①ある行に、0でない数をかける\\②ある行から、ほかの行の何倍かをひく\\③2つの行を入れかえる} & 7-5 \\
最小二乗法 & 正規方程式 $A^{\mathsf T}A\vec x=A^{\mathsf T}\vec b$（ずれの矢印が、列のつくる面に直角） & 7-5 \\
\end{longtblr}\par}

\bksubsection{5}{回転と座標の変換の公式}
{\small
\begin{longtblr}{colspec={Q[l,m,wd=26mm] Q[l,m,wd=86mm] Q[c,m,wd=12mm]}, rowhead=1, row{1}={bg=bkVol!14, font=\sffamily\bfseries\footnotesize}, hlines={0.4pt,bkLine}, rowsep=1.6pt}
公式・きまり & 式 & 章 \\
回転行列 & $R(\theta)=\begin{pmatrix}\cos\theta&-\sin\theta\\\sin\theta&\cos\theta\end{pmatrix}$（原点のまわりに左回りに $\theta$）。1列め＝$\vec e_x$ の行き先 & 7-6 \\
90°回す & $R_{90}=\begin{pmatrix}0&-1\\1&0\end{pmatrix}$、$(x,\ y)\to(-y,\ x)$。$R_{90}R_{90}=-I$ & 7-6 \\
つなぐ & $R(\alpha)R(\beta)=R(\alpha+\beta)$（中に加法定理）。$R(30°)R(45°)=R(75°)$ & 7-6 \\
もどす & $R(\theta)^{-1}=R(\theta)^{\mathsf T}=R(-\theta)$ & 7-6 \\
軸を回して見る & 軸を $\theta$ 回して見ると $\vec x\,'=R(\theta)^{\mathsf T}\vec x$（点は逆に回ったように見える） & 7-6 \\
2本の腕 & $\vec p=R(\theta_1)\begin{pmatrix}L_1\\0\end{pmatrix}+R(\theta_1)R(\theta_2)\begin{pmatrix}L_2\\0\end{pmatrix}$ & 7-6 \\
姿勢の角 & $R=R_z(\psi)R_y(\theta)R_x(\phi)$（ヨー $\psi$・ピッチ $\theta$・ロール $\phi$）。右から先にはたらく & 7-7 \\
機首の向き & $R$ の1列め $(\cos\psi\cos\theta,\ \sin\psi\cos\theta,\ -\sin\theta)$ & 7-7 \\
逆の回転 & $R^{-1}=R^{\mathsf T}=R_x(-\phi)R_y(-\theta)R_z(-\psi)$（順番も符号も逆） & 7-7 \\
同次変換行列 & $A(\theta,\ L)=\begin{pmatrix}\cos\theta&-\sin\theta&L\cos\theta\\\sin\theta&\cos\theta&L\sin\theta\\0&0&1\end{pmatrix}$。左上が回転、右の列が移動 & 7-7 \\
腕をつなぐ & $T=A(\theta_1,\ L_1)A(\theta_2,\ L_2)$（根もとが左）。$T$ の右の列が手の位置 & 7-7 \\
3次元 & 4×4 の $\begin{pmatrix}R&\vec p\\0\quad0\quad0&1\end{pmatrix}$（左上の $R$ は3×3、右上の $\vec p$ はたての3つ） & 7-7 \\
座標の変換 & $\vec p_{\text{地面}}=R\,\vec p_{\text{機体}}+\vec d$、\ $\vec p_{\text{機体}}=R^{\mathsf T}(\vec p_{\text{地面}}-\vec d)$ & 7-7 \\
\end{longtblr}\par}

\vspace{2mm}
\noindent\begin{minipage}{\linewidth}
\noindent{\sffamily\bfseries\small\color{\bkTT}3本の軸のまわりの回転}（7-7。回らない軸の行と列は「1と0」）\par\vspace{1mm}
\begin{bkcenter}
\begin{tblr}{colspec={Q[l,m,wd=44mm] Q[c,m,wd=62mm]}, column{1}={font=\sffamily\bfseries\footnotesize}, hlines={0.4pt,bkLine}, vlines={0.4pt,bkLine}, rowsep=2pt, cells={font=\small}}
{$x$ 軸のまわりに $\phi$ 回す\\（ロール）} & $R_x(\phi)=\begin{pmatrix}1&0&0\\0&\cos\phi&-\sin\phi\\0&\sin\phi&\cos\phi\end{pmatrix}$ \\
{$y$ 軸のまわりに $\theta$ 回す\\（ピッチ）} & $R_y(\theta)=\begin{pmatrix}\cos\theta&0&\sin\theta\\0&1&0\\-\sin\theta&0&\cos\theta\end{pmatrix}$ \\
{$z$ 軸のまわりに $\psi$ 回す\\（ヨー）} & $R_z(\psi)=\begin{pmatrix}\cos\psi&-\sin\psi&0\\\sin\psi&\cos\psi&0\\0&0&1\end{pmatrix}$ \\
\end{tblr}
\end{bkcenter}
{\footnotesize $R_y$ だけ、$\sin$ の符号が逆に見えます。$y$ 軸のまわりの＋は「$z$ を $x$ のほうへ」回すからです。$\theta=90°$ を入れ、3列めが $(1,\ 0,\ 0)$ になるかで確かめます（7-7）。}
\end{minipage}

\bksubsection{6}{固有値・対角化・くり返しの公式}
{\small
\begin{longtblr}{colspec={Q[l,m,wd=26mm] Q[l,m,wd=86mm] Q[c,m,wd=12mm]}, rowhead=1, row{1}={bg=bkVol!14, font=\sffamily\bfseries\footnotesize}, hlines={0.4pt,bkLine}, rowsep=1.6pt}
公式・きまり & 式 & 章 \\
固有値・固有ベクトル & $A\vec v=\lambda\vec v$（$\vec v\neq\vec 0$）。かけても向きの変わらない矢印 $\vec v$ と、その倍率 $\lambda$ & 7-8 \\
特性方程式 & $\det(A-\lambda I)=0$。2×2 なら $\lambda$ の2次方程式 & 7-8 \\
求め方 & {①対角成分から $\lambda$ をひいた $A-\lambda I$ を作る\\②$\det(A-\lambda I)=0$ を解いて、固有値を出す\\③固有値を1つずつ $(A-\lambda I)\vec v=\vec 0$ に入れ、$x$ と $y$ の関係を出す\\④簡単な $\vec v$ を答え、$A\vec v=\lambda\vec v$ をかけて確かめる} & 7-8 \\
確かめの式 & $\lambda_1+\lambda_2=a+d$、\ $\lambda_1\lambda_2=\det A$ & 7-8 \\
振動モード & $\ddot{\vec x}=B\vec x$ では、形は固有ベクトル、固有値は $-\omega^2$。$\omega=\sqrt{-\lambda}$ & 7-8 \\
6-12 の $s$ & $\dot{\vec x}=A\vec x$ に $\vec x=\vec ve^{st}$ を入れると $A\vec v=s\vec v$。$s$ は固有値 & 7-8 \\
対称な行列 & 固有値は実数、固有ベクトルは直交する & 7-8 \\
軸を変えて見る & $\vec x=P\vec q$、\ $\vec q=P^{-1}\vec x$（$P$ は固有ベクトルを列に並べた行列） & 7-9 \\
対角化 & $P^{-1}AP=D$、\ $A=PDP^{-1}$（$D$ は固有値を対角に並べた行列）。固有ベクトルが2本あるときだけ & 7-9 \\
対称な行列の対角化 & 長さを1にした固有ベクトルの回転行列 $R$ で $R^{\mathsf T}AR=D$ & 7-9 \\
何回もかける & $A^n=PD^nP^{-1}$。$D^n$ は対角成分を $n$ 乗するだけ & {7-9\\7-10} \\
くり返し & $\vec x_{n+1}=A\vec x_n$ なら $\vec x_n=A^n\vec x_0$。$A^n\vec v=\lambda^n\vec v$ & 7-10 \\
等比級数 & $1+x+x^2+\cdots=\dfrac{1}{1-x}$（$|x|<1$ のときだけ） & 7-10 \\
複素数の固有値 & $\lambda=re^{j\theta}$ なら $\lambda^n=r^ne^{jn\theta}$（大きさ $r^n$、回る角 $n\theta$） & 7-10 \\
オイラー法 & 1歩の行列 $\vec x_{n+1}=(I+A\,\Delta t)\,\vec x_n$。固有値は $1+\lambda\,\Delta t$（ほんとうは $e^{\lambda\Delta t}$）。こわれない条件 $|1+\lambda\,\Delta t|<1$ & 7-10 \\
\end{longtblr}\par}

\vspace{2mm}
\noindent\begin{minipage}{\linewidth}
\noindent{\sffamily\bfseries\small\color{\bkTT}固有値で、動きがわかる}（7-8・7-10）\par\vspace{1mm}
\begin{bkcenter}
\begin{tblr}{colspec={Q[l,m,wd=30mm] Q[l,m,wd=31mm] Q[l,m,wd=30mm] Q[l,m,wd=32mm]}, row{1}={bg=bkVol!14, font=\sffamily\bfseries\footnotesize}, row{2}={font=\sffamily\bfseries\footnotesize}, cell{2}{1}={c=2}{c}, cell{2}{3}={c=2}{c}, vline{3}={0.8pt,bkVol}, hlines={0.4pt,bkLine}, rowsep=1.4pt, cells={font=\small}}
\SetCell[c=2]{c} 微分方程式 $\dot{\vec x}=A\vec x$（7-8） & & \SetCell[c=2]{c} くり返し $\vec x_{n+1}=A\vec x_n$（7-10） & \\
固有値の実部で決まる & & 固有値の大きさ $|\lambda|$ で決まる & \\
負の実数 & ゆれずに小さくなる & $|\lambda|<1$ & 0に近づく（消える） \\
0 & 変わらない & $\lambda=1$ & その向きの部分が残る \\
正の実数 & どんどん大きくなる & $|\lambda|>1$ & ふえていく \\
複素数（実部が負） & ゆれながら小さくなる & 負の実数 & 向きが毎回入れかわる \\
複素数（実部が0） & 同じ大きさでゆれ続ける & 複素数 & 回りながら（大きさは $|\lambda|^n$） \\
\end{tblr}
\end{bkcenter}
{\footnotesize 何回もかけると、大きさ $|\lambda|$ がいちばん大きい固有値の固有ベクトルの向きにそろいます（べき乗法。7-8・7-10）。第9巻の制御では、固有値が全部「左半分」（実部が負）にあるように機械を作ります。}
\end{minipage}

\bksubsection{7}{表計算・Python の使い方}
{\small 表計算の行列は A2 から（1行めは見出し）、ベクトルはそのすぐ右の列、答えはさらにその右の列に置きます（7-1 の型）。}\par\vspace{1mm}
{\small
\begin{longtblr}{colspec={Q[l,m,wd=27mm] Q[l,m,wd=55mm] Q[l,m,wd=33mm] Q[c,m,wd=10mm]}, rowhead=1, row{1}={bg=bkVol!14, font=\sffamily\bfseries\footnotesize}, column{2}={font=\footnotesize}, hlines={0.4pt,bkLine}, rowsep=1.6pt, colsep=3pt}
すること & 書き方（例） & 意味・答え & 章 \\
{行列×ベクトル\\（表計算）} & \texttt{=MMULT(A2:B3,C2:C3)} & {D2:D3 に出る\\（左の範囲、右の範囲の順）} & 7-1 \\
{行列×行列\\（表計算）} & \texttt{=MMULT(A2:B3,C2:D3)} & E2:F3 に出る & 7-1 \\
転置 & \texttt{=TRANSPOSE(A2:B3)} & 行と列を入れかえる & 7-1 \\
行列式 & \texttt{=MDETERM(A2:B3)} & 主役の行列なら3 & 7-4 \\
{逆行列で解く\\（表計算）} & \texttt{=MMULT(MINVERSE(A2:B3),C2:C3)} & 3-6 の回路なら2、0.5 & 7-5 \\
最小二乗法 & \texttt{=LINEST(B2:B8,A2:A8)} & {1.4893、0.0321\\（$y$ の範囲が先）} & 7-5 \\
回転行列 & {\texttt{=COS(RADIANS(F2))}\\\texttt{=-SIN(RADIANS(F2))}} & F2 に角（度） & 7-6 \\
$P^{-1}AP$ & {\texttt{=MMULT(MMULT(MINVERSE(A2:B3),}\\\texttt{\ \ C2:D3),A2:B3)}} & {$P$ は A2:B3、\\$A$ は C2:D3} & 7-9 \\
{行列を作る\\（Python）} & {\texttt{import numpy as np}\\\texttt{A = np.array([[2, 1], [1, 2]])}} & 1行ずつ \texttt{[ ]} でくくる & 7-1 \\
積 & \texttt{A @ x}、\texttt{A @ B} & \texttt{*} は成分ごとの積 & 7-1 \\
転置・単位行列 & \texttt{A.T}、\texttt{np.eye(2)} & & 7-1 \\
行列式 & \texttt{np.linalg.det(A)} & 3（\texttt{2.999…} と出る） & 7-4 \\
逆行列 & \texttt{np.linalg.inv(A)} & & 7-5 \\
連立方程式 & \texttt{np.linalg.solve(A, b)} & $A\vec x=\vec b$ の $\vec x$ & 7-5 \\
最小二乗法 & \texttt{np.linalg.lstsq(A, y, rcond=None)[0]} & $a$、$b$ & 7-5 \\
度をラジアンに & \texttt{np.radians(30)} & \texttt{np.cos} はラジアン & 7-6 \\
固有値と固有ベクトル & \texttt{w, V = np.linalg.eig(A)} & {\texttt{w} に固有値\\\texttt{V} の列に固有ベクトル} & 7-8 \\
固有値だけ・大きさ & \texttt{abs(np.linalg.eigvals(G))} & 大きさ $|\lambda|$ & 7-10 \\
$n$ 乗 & \texttt{np.linalg.matrix\_power(A, 5)} & \texttt{A ** 5} は成分ごと & 7-10 \\
\end{longtblr}\par}
\begin{itemize}[leftmargin=1.4\zw, itemsep=2pt, topsep=4pt, label={\color{\bkTT}●}]
\item 表計算の確定のしかた：新しい Excel（Microsoft 365）や Google スプレッドシートでは、答えの左上のセルに式を入れて Enter を押すと、答えが自動で広がります。広がらないソフト（古い Excel や LibreOffice）では、先に答えのマス全部をえらんでから式を入れ、Ctrl＋Shift＋Enter で確定します（7-1）。
\item 大きさが合わないときや、範囲に空いたセルがあるときは \texttt{\#VALUE!}、行列式が0で逆行列がないときは \texttt{\#NUM!} と出ます。Python では \texttt{LinAlgError: Singular matrix} と出ます（7-1・7-5）。
\item \texttt{2.9999999999999996} は、2進法で計算したときのごくわずかなずれです。行列式は3です。\texttt{round(np.linalg.det(A), 6)} で \texttt{3.0} になります（7-4）。
\item \texttt{3.+0.j} は「虚部が0の複素数」で、3のことです。numpy 2.0 より新しい版では、固有値はいつも複素数の形で出ます（7-8）。
\item numpy の固有ベクトルは\textbf{列}（たて）に並び、長さが1にそろえてあります。$0.70710678=\dfrac{1}{\sqrt2}$ です。向きが反対でも、同じ直線の上なら同じ固有ベクトルです（7-8）。
\item ベクトルは \texttt{np.array([120, 60])} と横に書き、表示も \texttt{[420 300]} と横1行ですが、中身は数学のたて書きの列ベクトルと同じです（7-1）。
\end{itemize}

\Needspace{8\baselineskip}
\bkbkhead{プログラムのひな形（7-5・7-8）}
{\small 上は 3-6 の回路の連立方程式（行列式と、\texttt{solve} の答え $I_1=\qty{2}{A}$、$I_2=\qty{0.5}{A}$）。下は主役の行列の固有値と固有ベクトル（\texttt{V} の1列めが固有値3の組）。左がプログラム、右が実行結果です。数を入れかえれば、ほかの行列にも使えます。}\par\vspace{1mm}
\noindent\begin{minipage}[t]{0.42\linewidth}\vspace{0pt}
\begin{bkcode}[Python]
import numpy as np
A = np.array([[4, 2],
              [2, 4]])
b = np.array([9, 6])
print(np.linalg.det(A))
x = np.linalg.solve(A, b)
print(x)
\end{bkcode}
\end{minipage}\hfill
\begin{minipage}[t]{0.55\linewidth}\vspace{0pt}
\begin{bkcodeout}[実行結果]
12.0
[2.  0.5]
\end{bkcodeout}
\end{minipage}\par\vspace{2mm}
\noindent\begin{minipage}[t]{0.42\linewidth}\vspace{0pt}
\begin{bkcode}[Python]
import numpy as np
A = np.array([[2, 1],
              [1, 2]])
w, V = np.linalg.eig(A)
print(w)
print(V)
\end{bkcode}
\end{minipage}\hfill
\begin{minipage}[t]{0.55\linewidth}\vspace{0pt}
\begin{bkcodeout}[実行結果]
[3.+0.j 1.+0.j]
[[ 0.70710678+0.j -0.70710678+0.j]
 [ 0.70710678+0.j  0.70710678+0.j]]
\end{bkcodeout}
\end{minipage}\par\vspace{2mm}

\bksubsection{8}{専門書の書き方との対応}
{\small 大学の本では、同じものを別の記号やことばで書きます。各章の「専門書をのぞいてみよう」で読みかえたものを集めました。}\par\vspace{1mm}
{\small
\begin{longtblr}{colspec={Q[l,m,wd=44mm] Q[l,m,wd=66mm] Q[c,m,wd=14mm]}, rowhead=1, row{1}={bg=bkVol!14, font=\sffamily\bfseries\footnotesize}, hlines={0.4pt,bkLine}, rowsep=1.2pt}
専門書の書き方 & この本の書き方・意味 & 章 \\
太字の $\symbf{A}$、$\symbf{x}$ & 細字の行列 $A$、矢印のベクトル $\vec x$ & {7-2\\7-11} \\
$A=(a_{ij})$ & 成分が $a_{ij}$ の行列 $A$ & 7-1 \\
${}^tA$ & 転置 $A^{\mathsf T}$ & {7-1\\7-6} \\
$E$ & 単位行列 $I$ & {7-1\\7-5} \\
$|A|$、\ $|A-\lambda E|$ & 行列式 $\det A$、\ $\det(A-\lambda I)$ & {7-4\\7-8} \\
写像、線形写像、一次変換 & 行列をかける変換（入れると出てくる機械） & {7-2\\7-6} \\
$\mathbb{R}^2$、\ $\mathbb{R}^n$ & 2つの数の組（平面の点）全部、$n$ 個の実数の組の全体 & {7-2\\7-11} \\
$n$ 次正方行列 & たて・横 $n$ 個ずつの行列 & 7-8 \\
零ベクトル & $\vec 0$ & 7-8 \\
標準基底 $\symbf{e}_1$、$\symbf{e}_2$ & 物差しの矢印 $\vec e_x$、$\vec e_y$ & 7-2 \\
固有方程式 & 特性方程式 $\det(A-\lambda I)=0$ & 7-8 \\
正則 & 逆行列がある & 7-5 \\
ガウスの消去法 & 掃き出し法 & 7-5 \\
主応力、主方向 & 応力の行列の固有値、固有ベクトルの向き & 7-8 \\
一次独立（2本のとき） & 平行でない & 7-9 \\
$\mathrm{diag}(\lambda_1,\ \lambda_2)$ & 固有値を対角に並べた対角行列 $D$ & 7-9 \\
直交行列 & 回転（や裏返し）の行列。$R^{\mathsf T}=R^{-1}$ & 7-9 \\
二端子対回路、縦続接続 & 入口2本・出口2本の箱、箱をくさりのようにつなぐこと & 7-3 \\
四端子定数 $A$、$B$、$C$、$D$（F 行列） & 箱の行列の4つの数（行列の名前ではない） & 7-3 \\
${}^{i-1}\symbf{T}_i$、\ ${}^0\symbf{T}_n$ & 1本の腕の $A(\theta_i,\ L_i)$、つないだ全体の $T$ & 7-7 \\
ZYX オイラー角 & $R_z(\psi)R_y(\theta)R_x(\phi)$ の順の姿勢の角 & 7-7 \\
差分方程式、前進オイラー法の $h$ & 漸化式、ふつうのオイラー法の $\Delta t$ & 7-10 \\
\end{longtblr}\par}

%% ---------------- チェック表 ----------------
\bkbacktitle{できるようになったことチェック表}
{\small 各章の「ふりかえり」を集めた表です。章を学び終えたら1回め、次の巻に進む前に2回めのチェックをしよう。できなかった所は、ページにもどって読みなおします。}\par\vspace{2mm}
%% 自動生成（tools/mkchecklist.py）。各章の「ふりかえり」の項目を集めたもの
\begingroup\footnotesize
\def\bkcb{\tikz[baseline=-0.5ex]\draw[\bkTF, line width=0.7pt, rounded corners=0.5pt] (0,-0.3em) rectangle (0.8em,0.5em);}
\begin{longtblr}{colspec={Q[c,m,wd=9mm] X[l,m] Q[c,m,wd=10mm] Q[c,m,wd=9mm] Q[c,m,wd=9mm]}, width=\linewidth, rowhead=1, colsep=3pt,
  row{1}={bg=bkVol!14, font=\sffamily\bfseries\scriptsize}, hlines={0.4pt,bkLine}, rowsep=0.7pt}
章 & できるようになったこと & ページ & 1回め & 2回め \\
{\sffamily\bfseries\color{\bkTT}7-1} & 数の表を行列として読み、成分 $a_{ij}$ を言える。和・差・実数倍を計算できる。 & \pageref{bk:7-1:gyoretsu} & \bkcb & \bkcb \\
 & 行列×ベクトル・行列×行列を、行と列をなぞって計算できる。 & \pageref{bk:7-1:kakeru} & \bkcb & \bkcb \\
 & 連立方程式を $A\vec x=\vec b$ の形に書き、表計算や Python で確かめられる。 & \pageref{bk:7-1:renritsu} & \bkcb & \bkcb \\
{\sffamily\bfseries\color{\bkTT}7-2} & 行列×ベクトルを、行と列をなぞって計算できる。 & \pageref{bk:7-2:nazoru} & \bkcb & \bkcb \\
 & 列（物差しの矢印の行き先）から、絵の変わり方を読みとり、行き先から行列を作れる。 & \pageref{bk:7-2:retsu} & \bkcb & \bkcb \\
 & ヤコビ行列や応力の行列を、「入れると出てくる機械」として計算できる。 & \pageref{bk:7-2:jacobi} & \bkcb & \bkcb \\
{\sffamily\bfseries\color{\bkTT}7-3} & 行列×行列を、行と列の内積で計算できる。 & \pageref{bk:7-3:naiseki} & \bkcb & \bkcb \\
 & 変換をつなぐと積になり、順番で結果が変わることを説明できる。 & \pageref{bk:7-3:junban} & \bkcb & \bkcb \\
 & 電気の箱をつないだ回路を、行列の積で計算できる。 & \pageref{bk:7-3:hako} & \bkcb & \bkcb \\
{\sffamily\bfseries\color{\bkTT}7-4} & 方眼で面積を数え、$ad-bc$ で行列式を計算できる。 & \pageref{bk:7-4:adbc} & \bkcb & \bkcb \\
 & 行列式のマイナス（裏返し）と0（平らにつぶれる）の意味がわかる。 & \pageref{bk:7-4:zero} & \bkcb & \bkcb \\
 & 3×3 の行列式をサラスの方法で計算し、工学の行列式を読める。 & \pageref{bk:7-4:sanjigen} & \bkcb & \bkcb \\
{\sffamily\bfseries\color{\bkTT}7-5} & 2×2 の逆行列を作って確かめ、$\vec x=A^{-1}\vec b$ で連立方程式を解ける。 & \pageref{bk:7-5:gyaku} & \bkcb & \bkcb \\
 & 掃き出し法で、3つの文字の連立方程式と逆行列を求められる。 & \pageref{bk:7-5:hakidashi} & \bkcb & \bkcb \\
 & 正規方程式 $A^{\mathsf T}A\vec x=A^{\mathsf T}\vec b$ で、最小二乗法の直線を求められる。 & \pageref{bk:7-5:saisho} & \bkcb & \bkcb \\
{\sffamily\bfseries\color{\bkTT}7-6} & 回転行列 $R(\theta)$ を作り、点や図形を回せる。 & \pageref{bk:7-6:tsukuru} & \bkcb & \bkcb \\
 & $R(\alpha)R(\beta)=R(\alpha+\beta)$ と $R(\theta)^{-1}=R(\theta)^{\mathsf T}$ を使える。 & \pageref{bk:7-6:tsunagu} & \bkcb & \bkcb \\
 & 軸を回して点や力を見たり、腕の手の位置を求めたりできる。 & \pageref{bk:7-6:jiku} & \bkcb & \bkcb \\
{\sffamily\bfseries\color{\bkTT}7-7} & 3本の軸のまわりの回転を $R_x$・$R_y$・$R_z$ で表し、点を回せる。 & \pageref{bk:7-7:rz} & \bkcb & \bkcb \\
 & ロール・ピッチ・ヨーから機首の向きを求め、順番で向きが変わることを説明できる。 & \pageref{bk:7-7:euler} & \bkcb & \bkcb \\
 & 同次変換行列で、腕の手の位置やカメラで見た物の位置を求められる。 & \pageref{bk:7-7:doji} & \bkcb & \bkcb \\
{\sffamily\bfseries\color{\bkTT}7-8} & 行列をかけても向きが変わらない矢印を、図と表計算でさがせる。 & \pageref{bk:7-8:sagasu} & \bkcb & \bkcb \\
 & $A\vec v$ を計算して固有ベクトルを確かめ、振動モードが固有ベクトルだと説明できる。 & \pageref{bk:7-8:daisha} & \bkcb & \bkcb \\
 & 特性方程式で固有値と固有ベクトルを求め、6-12 の $s$ と同じものだと確かめられる。 & \pageref{bk:7-8:tokusei} & \bkcb & \bkcb \\
{\sffamily\bfseries\color{\bkTT}7-9} & 固有ベクトルの物差しで見ると「のびるだけ」になることを、図で説明できる。 & \pageref{bk:7-9:naname} & \bkcb & \bkcb \\
 & $P$ と $D$ を作り、$P^{-1}AP=D$ を計算で確かめられる。 & \pageref{bk:7-9:jiku} & \bkcb & \bkcb \\
 & 対角化で $A^n$ を計算したり、台車の式をモードの座標でほどいたりできる。 & \pageref{bk:7-9:hodoku} & \bkcb & \bkcb \\
{\sffamily\bfseries\color{\bkTT}7-10} & 同じ行列を何回もかけて、点の向かう先を表計算と図で調べられる。 & \pageref{bk:7-10:kurikaeshi} & \bkcb & \bkcb \\
 & 倍々のくり返しと、等比級数の合計を計算できる。 & \pageref{bk:7-10:kyusu} & \bkcb & \bkcb \\
 & 固有値の大きさで行き先を見分け、ばねがふくらむわけを説明できる。 & \pageref{bk:7-10:bane} & \bkcb & \bkcb \\
{\sffamily\bfseries\color{\bkTT}7-11} & ロボットの腕を、ヤコビ行列の逆行列で目標へ近づけるしくみを説明できる。 & \pageref{bk:7-11:gyaku} & \bkcb & \bkcb \\
 & ばねの網と抵抗の網、箱をつないだ回路を、行列で書ける。 & \pageref{bk:7-11:ami} & \bkcb & \bkcb \\
 & 長い電線が $\cosh$・$\sinh$ の行列になるわけを、箱のくさりとして読める。 & \pageref{bk:7-11:bunpu} & \bkcb & \bkcb \\
\end{longtblr}
\endgroup


%% ---------------- つぎの巻へ（第8巻の予告と、この巻の道具がこの先どこで使われるか） ----------------
\bkbacktitle{つぎの巻へ}
\newcommand\bknextpill[1]{\bk@pill{#1}}
\def\bknextrow{1}
\newcommand\bknextcard[4]{%
  \begin{tcolorbox}[enhanced, nobeforeafter, box align=top, width=0.485\linewidth, colframe=bkVol8, colback=white, boxrule=0.9pt, arc=2.5mm,
    left=2.5mm, right=2.5mm, top=1.5mm, bottom=2mm, equal height group=bknext\bknextrow]
  \raggedright\bknextpill{#1}\enspace{\sffamily\bfseries\small #2}\par\vspace{1mm}
  {\sffamily\bfseries\footnotesize\color{bkDeep8}#3}\par\vspace{1mm}
  {\footnotesize\setlength{\parindent}{0pt}{\sffamily\bfseries\color{bkInk}使う道具}\enspace #4}
  \end{tcolorbox}}
たくさんの数を、表のまま計算する道具がそろいました。第8巻「空間に広がる量」では、空間のどの点にも矢印がある量を、微分と積分で読みときます。7-2 のヤコビ行列と 7-4 の行列式が、その入口です。\par\vspace{2.5mm}
\noindent\bknextcard{8-3}{全微分：少し動いたときの変化}{ななめに少し進むと、高さはどれだけ変わる？}{ヤコビ行列 $J$ と行列×ベクトル（7-2）。}\hfill
\bknextcard{8-4}{いちばん急な向き：勾配 grad}{ボールが転がる向きは、等高線と直角}{行と列の内積（7-3）。内積が0なら直角。}\par\vspace{1.2mm}
\def\bknextrow{2}\noindent\bknextcard{8-5}{ベクトル場}{空間の点ごとに、矢印を1本ずつ}{3次元の回転と座標の変換（7-7）。}\hfill
\bknextcard{8-6}{面や立体でたし合わせる：重積分}{厚さがちがう板の重さを求める}{行列式は面積の倍率（7-4）。ヤコビ行列の行列式も使う。}\par\vspace{1.2mm}
\def\bknextrow{3}\noindent\bknextcard{8-10}{うず：回転 rot}{小さな風車を置くと、回るかな？}{外積（4-14）と、3×3 の行列式（7-4）。}\hfill
\bknextcard{8-13}{熱と波の伝わり方：偏微分方程式}{金属の棒に熱が広がるようすを、表計算で再現する}{1歩の行列のくり返しと、こわれない条件（7-10）。}\par

\clearpage
\bkbkhead{この巻の道具は、この先でこう使う}
{\small 第7巻の章と、その道具を使う第8巻から先の章です（章のはじめの「道しるべ」の「この章の先に」から集めました）。}\par\vspace{3mm}
%% 自動生成（tools/mknext.py）。各章の道しるべ「この章の先に」から、後ろの巻の章を集めたもの
\begin{longtblr}{colspec={Q[l,m,wd=58mm] X[l,m]}, width=\linewidth, rowhead=1, row{1}={bg=bkVol!14, font=\sffamily\bfseries\footnotesize},
  hlines={0.4pt,bkLine}, rowsep=1.5pt, cells={font=\footnotesize}}
この巻の章 & この道具を使う章 \\
{\bknextpill{7-2}\ 行列×ベクトル：形を変える機械} & {\bknextpill{8-3}\ 全微分：少し動いたときの変化} \\
{\bknextpill{7-5}\ 逆行列で連立方程式を解く} & {\bknextpill{9-10}\ 連立微分方程式を行列で：状態方程式} \\
{\bknextpill{7-7}\ 3次元の回転と座標の変換} & {\bknextpill{8-5}\ ベクトル場} \\
{\bknextpill{7-8}\ 固有値と固有ベクトル：向きの変わらない方向} & {\bknextpill{9-10}\ 連立微分方程式を行列で：状態方程式\\ \bknextpill{9-11}\ 行列の指数関数と固有値} \\
{\bknextpill{7-9}\ 対角化：まっすぐな方向に分けて考える} & {\bknextpill{9-10}\ 連立微分方程式を行列で：状態方程式\\ \bknextpill{9-11}\ 行列の指数関数と固有値} \\
{\bknextpill{7-10}\ 行列のくり返し：漸化式とオイラー法を行列で} & {\bknextpill{9-10}\ 連立微分方程式を行列で：状態方程式\\ \bknextpill{9-11}\ 行列の指数関数と固有値} \\
{\bknextpill{7-11}\ 行列の使いみち} & {\bknextpill{8-3}\ 全微分：少し動いたときの変化\\ \bknextpill{9-10}\ 連立微分方程式を行列で：状態方程式\\ \bknextpill{9-11}\ 行列の指数関数と固有値\\ \bknextpill{9-12}\ フーリエとラプラスの使いみち} \\
\end{longtblr}

\vspace{4mm}
\bkfukidashi{delta}{第7巻の行列は、第8巻・第9巻・物理シリーズでずっと使うよ。自信のない章を見なおしてから、第8巻へ進もう！}

%% ---------------- さくいん ----------------
\bkbacktitle{さくいん}
{\small 数字はページです。その語がはじめて出てくるページや、くわしく説明しているページをのせています。}\par\vspace{2mm}
\begin{multicols}{2}\small\setlength{\parskip}{0pt}
\directlua{bk.printidx()}
\end{multicols}

%% ---------------- 奥付 ----------------
\bkpadfour
\global\bk@decofalse\bkfolio{false}
\null\vfill
\noindent\begin{minipage}{\linewidth}\small
{\sffamily\bfseries\Large\color{bkDeep}物理を理解するための数学　第7巻}\par\vspace{1mm}
{\sffamily\bfseries\large まとめて計算する}\par\vspace{3mm}
{\color{bkVol}\hrule height 1pt}\vspace{3mm}
\begin{tabular}{@{}l@{\hspace{6mm}}l@{}}
編者 & 「物理を理解するための数学」編集委員会 \\
 & 委員長　小学校長（まとめ役） \\
 & 委員　数学・電気工学・機械工学・教育学・量子力学・建築工学・\\
 & \phantom{委員　}半導体工学の大学教授、幼稚園教諭、塾講師、\\
 & \phantom{委員　}特別支援学校教諭、心理士、科学解説者、小学校教諭 \\
デザイン & ユニバーサルデザイン担当委員 \\
挿絵 & イラスト担当委員 \\
版 & 第1版（第14回編集委員会）　2026年10月 \\
\end{tabular}\par\vspace{3mm}
{\color{bkVol}\hrule height 1pt}\vspace{2mm}
{\footnotesize この本の問題の答えは、巻末の「こたえ」にあります。くわしい解き方は、別冊の問題集と指導書にのせます。}
\end{minipage}
\makeatother
